% =====================================================================
%  PRIMER PART 3 — RESULTS WITH FULL PEDAGOGICAL PROOFS
%  The answers to the questions raised in Part II, proved in detail
%  for a strong undergraduate who has done Part I (robustness,
%  Lipschitz/eta-L certificates, margins/convex hulls, implicit bias,
%  NTK/lazy training, Fourier/power-spectrum/bispectrum invariants).
%
%  This is a FRAGMENT: no \documentclass, no \begin{document}.
%  The assembler (primer.tex) supplies the preamble, which already
%  defines: definition, theorem, lemma, proposition, corollary,
%  conjecture, example, remark; and the macros
%  \R \Z \E \sign \Lip \one \conv \dist \Orb \Vinv \Vac \Piinv
%  \fdc \Sep .
%
%  Proofs are EXPANDED line-by-line for a strong undergraduate.
%  Citation keys used appear in the comment block at the END of file
%  (new keys are flagged).
% =====================================================================

\section{Results: what survives an invariance, and what it costs}
\label{sec:results}

Part~II asked five questions. Each came from a real friction point: the old project's monotone
``more consistency, less robustness'' story broke and could not be explained; the base papers
openly contradict each other on whether invariance helps or hurts; and no prior result ties a
\emph{translation} invariance to an $\ell_2$ robust radius. This part answers those questions with
theorems and full proofs. Two honesty rules run through the whole part, and we state them once,
plainly, so no proof reads as a bigger claim than it is.

\begin{remark}[Calibration: what is new here and what is borrowed]
\label{rem:calibration}
Most of the \emph{machinery} below is classical. Group averaging, fixed subspaces, the
Lipschitz--margin certificate, power-spectrum and bispectrum invariants, the implicit-bias
characterization of gradient flow, and the lazy/NTK picture are all established tools, taught in
Part~I and cited there. What this part contributes is narrow and specific:
(1)~an \emph{exact projection-margin accounting} of what discriminative signal a given invariance
keeps (Theorem~\ref{thm:fg-margin});
(2)~a \emph{two-sided bracket} on the robust radius (a pair of bounds, lower and upper, that trap the
true radius from both sides) that completes the missing upper arm of the known certificate chain
(Theorem~\ref{thm:sandwich});
(3)~a \emph{negative} optimization result, that on \emph{orbit-closed} data (data where every shift
of a training point is itself a training point of the same label) imposing invariance is margin-free
(Theorem~\ref{thm:selfsym});
(4)~a \emph{positive} optimization result in the \emph{lazy} (wide-network, fixed-kernel) regime,
that invariance lowers the Lipschitz constant at equal margin (Theorem~\ref{thm:lazy}).
We never call $\eta/L$ a ``new robustness quantity'': it is a standard certificate, and our angle
is only to compute it \emph{after} identifying the invariant feature space and then to track which
of its two parts (the margin $\eta$ or the Lipschitz term $L$) a given architecture moves.
\end{remark}

\begin{remark}[Two plain-language terms we will need]
\label{rem:terms}
A \emph{converse} to an inequality $A \le B$ is a bound in the other direction, $B \le C\,A$ for
some constant $C$: it says $A$ is not just a lower estimate of $B$ but pins $B$ down to within a
factor. A certificate $r_2 \ge \eta/L$ becomes \emph{useful as a predictor} only if it has a
converse, otherwise the true radius could be arbitrarily larger than the certificate.
A function $g$ is \emph{co-Lipschitz} (or lower-Lipschitz) with constant $\alpha$ along a path if it
\emph{cannot} change too slowly: $|g(z)-g(x)| \ge \alpha\,\|z-x\|_2$. Ordinary Lipschitz continuity
($|g(z)-g(x)| \le L\,\|z-x\|_2$) caps how fast $g$ changes; the co-Lipschitz bound puts a floor
under it. The certificate's missing converse is exactly a co-Lipschitz statement, as we will see.
\end{remark}

\begin{remark}[Notation: one margin, several names]
\label{rem:margin-notation}
The same conceptual quantity --- the \emph{margin}, ``how decisively is a point classified'' ---
wears different symbols depending on the setting, and we flag this once so it never confuses.
In Part~I's robustness section the \emph{logit} margin (top logit minus best competitor) was
$M_{F,x}$; in Part~I's learning section the \emph{geometric} (SVM) margin, a distance to the
boundary, was $\gamma$. In this part we write $\eta$ for the \emph{feature margin} of an invariant
score (the gap the surviving signal supports) and $\mu = y\,g(x)$ for the signed margin used inside
the radius bracket; for a max-margin separator these coincide, $\mu = \eta$. Likewise $r_2$ here is
the $\ell_2$ robust radius written $r(x)$ in Part~I. Finally, Part~III's Fourier proofs use the
\emph{unitary} DFT $\hat x_k = d^{-1/2}\sum_j x_j e^{-2\pi i jk/d}$ (so Parseval is an exact isometry),
whereas Part~I fixed the \emph{unnormalised} forward DFT (so Parseval carried a factor of $N$); both
use the same minus-sign root of unity $\omega = e^{-2\pi i/d}$.
\end{remark}


% =====================================================================
\subsection{Q2 answered: the invariant linear margin is a projected convex-hull distance}
\label{sub:fg-margin}

\paragraph{The question this answers.}
A shift-invariant linear classifier must give every circular shift of a signal the same score. So
it cannot use any feature that a shift can change. \emph{Which} discriminative signal is left? Part~I
showed (for cyclic shifts) that the fixed subspace is the one-dimensional line spanned by the
all-ones vector, so an invariant linear score sees only the mean (DC) component. The clean general
statement is not special to cyclic shifts or to the DC line: for \emph{any} finite group of
orthogonal symmetries, the surviving signal is the data \emph{projected onto the fixed subspace},
and the best achievable invariant margin is exactly half the distance between the two projected
convex hulls. This is the structural anchor for everything that follows.

\paragraph{Status.}
This is a clean generalization, not a new phenomenon. \citet{ge2021shift} prove the cyclic/DC special
case (an invariant linear classifier is governed by the DC gap), and group averaging onto a fixed
subspace is standard representation theory. The modest contribution is the general finite-group
statement with the convex-hull normalization of the margin; the cyclic case recovers Ge, including
the $2/\sqrt{d}$ collapse, as a corollary. We do not claim the cyclic effect as new.

We use one fact from Part~I about two finite point sets $A,B$ in $\R^m$: the largest margin of a
hyperplane that separates them, where the margin is the smallest distance from any point to the
hyperplane, equals $\tfrac12\dist(\conv A,\conv B)$, the half-distance between their convex hulls.
This is the geometric (convex-hull) form of the hard-margin support vector machine
\citep{bennett2000,burges1998}. We will reuse it verbatim, so we restate it as the lemma we lean on.

\begin{lemma}[Hard-margin = half convex-hull distance]
\label{lem:hull}
For finite sets $A,B \subset \R^m$ with $\conv A \cap \conv B = \varnothing$, the maximum over unit
vectors $n$ and biases $b$ of the separation margin
\[
  \min\big(\,\min_{a\in A}(n^\top a + b),\ \min_{b'\in B}(-n^\top b' - b)\,\big)
\]
equals $\tfrac12\dist(\conv A,\conv B)$, and is attained by the unit normal pointing along the
segment joining the two closest hull points, with the hyperplane placed at its midpoint.
\end{lemma}

\begin{theorem}[Invariant linear margin for a finite orthogonal group]
\label{thm:fg-margin}
Let a finite group $G$ act on $\R^d$ by orthogonal matrices $\{R_g\}_{g\in G}$. Define the fixed
subspace and the group-average (Reynolds) projector
\[
V^G := \{v : R_g v = v \ \text{for all } g\in G\}, \qquad
P_G := \frac{1}{|G|}\sum_{g\in G} R_g .
\]
A linear score $h(x) = w^\top x + b$ is exactly invariant, $h(R_g x) = h(x)$ for all $x$ and $g$, if
and only if $w \in V^G$. For finite labelled sets $X_+, X_- \subset \R^d$, the largest $\ell_2$
margin achievable by an exactly invariant linear classifier is
\[
\gamma_{\mathrm{inv}}
= \tfrac12 \dist\big(\conv(P_G X_+),\ \conv(P_G X_-)\big).
\]
\end{theorem}

\begin{proof}
\emph{Step 1 (which normals are invariant).}
Fix $g$. The identity $h(R_g x) = h(x)$ for \emph{all} $x$ means $w^\top R_g x + b = w^\top x + b$
for all $x$, i.e. $w^\top R_g x = w^\top x$ for all $x$. Subtracting, $(R_g^\top w - w)^\top x = 0$
for all $x \in \R^d$. A vector orthogonal to all of $\R^d$ is zero, so $R_g^\top w = w$.
\emph{(orthogonality of the action)} Because $R_g$ is orthogonal, $R_g^\top = R_g^{-1} = R_{g^{-1}}$,
and as $g$ ranges over the group so does $g^{-1}$; hence $R_g^\top w = w$ for all $g$ is the same
condition as $R_g w = w$ for all $g$, that is, $w \in V^G$. The converse is the same computation read
backwards: if $w \in V^G$ then $R_g^\top w = w$ gives $w^\top R_g x = w^\top x$.

\emph{Step 2 ($P_G$ is the orthogonal projector onto $V^G$).}
First $P_G$ is idempotent: \emph{(group closure)}
\[
P_G^2 = \frac{1}{|G|^2}\sum_{g,h\in G} R_g R_h
       = \frac{1}{|G|^2}\sum_{g,h} R_{gh}
       = \frac{1}{|G|^2}\sum_{h}\Big(\sum_g R_{gh}\Big)
       = \frac{1}{|G|^2}\sum_{h} |G|\,P_G = P_G,
\]
where for each fixed $h$ the inner sum over $g$ runs over all group elements $gh$, hence equals
$|G|\,P_G$. Next $P_G$ is symmetric: \emph{(orthogonality again)}
$P_G^\top = \frac{1}{|G|}\sum_g R_g^\top = \frac{1}{|G|}\sum_g R_{g^{-1}} = P_G$. A symmetric
idempotent is an orthogonal projector. Its range is $V^G$: if $v\in V^G$ then $P_G v = \frac1{|G|}
\sum_g R_g v = \frac1{|G|}\sum_g v = v$, so $V^G \subseteq \mathrm{range}(P_G)$; and for any $x$,
$R_h(P_G x) = \frac1{|G|}\sum_g R_{hg} x = P_G x$ by the same reindexing, so
$\mathrm{range}(P_G)\subseteq V^G$.

\emph{Step 3 (invariant classification = classification of the projection).}
Take any invariant $w \in V^G$. By Step~2, $w = P_G w$, and since $P_G$ is symmetric,
\[
w^\top x = (P_G w)^\top x = w^\top P_G x .
\]
So evaluating the invariant score at $x$ is the same as evaluating the \emph{same} normal $w$ at the
projected point $P_G x$, with the same $\|w\|_2$. Therefore an invariant linear classifier on
$X_+ \cup X_-$ is exactly an ordinary linear classifier on the projected sets $P_G X_+$ and
$P_G X_-$, and the signed margins match point by point. Conversely any linear classifier in the
projected space lifts (take its normal, which already lies in $V^G$) to an invariant one.

\emph{Step 4 (apply the hard-margin lemma).}
By Step~3, maximizing the invariant margin over original-space data is the same problem as
maximizing the ordinary linear margin over the projected data $P_G X_+$, $P_G X_-$.
Lemma~\ref{lem:hull} with $A = P_G X_+$ and $B = P_G X_-$ gives that this maximum equals
$\tfrac12\dist(\conv(P_G X_+),\conv(P_G X_-))$, attained by the closest-points normal and the
midpoint bias. This is the claim.
\end{proof}

\begin{remark}[Why the right object is a distance, not a one-sided gap]
A natural but wrong guess writes the invariant margin as a labelled gap such as
$\tfrac12(\min_{X_+}\fdc - \max_{X_-}\fdc)_+$, which is only correct when the positive class happens
to sit ``above'' the negative one along the fixed direction. The margin is symmetric in the labels:
it is the \emph{distance} between the two projected convex hulls, which is orientation-free and is
zero exactly when they overlap. This is the symmetric form proved above.
\end{remark}

\paragraph{Recovering Ge's $2/\sqrt d$ collapse.}
For cyclic shifts on $\R^d$ the fixed subspace is the single line $\mathrm{span}\{\one\}$ and
$P_G = \Piinv = ee^\top$ with $e = d^{-1/2}\one$ (Part~I). Then $P_G x = (e^\top x)e = \fdc(x)\,e$,
so the projected data lie on one line and are identified isometrically with the scalar DC values
$\fdc(x)$. The two projected hulls become two intervals $I_+, I_-$ on that line, and
Theorem~\ref{thm:fg-margin} reduces to
\[
\gamma_{\mathrm{inv}} = \tfrac12\dist(I_+, I_-).
\]
Take Ge's black/white dot example: $x_+ = e_j$ and $x_- = -e_j$ (a single bright pixel versus a
single dark pixel). Their DC values are $\fdc(\pm e_j) = \pm d^{-1/2}$, so the interval distance is
$2/\sqrt d$ and the invariant margin is $\gamma_{\mathrm{inv}} = 1/\sqrt d$. Yet \emph{without} the
invariance constraint a linear classifier may use the coordinate $x_j$ directly, separating
$\pm e_j$ with margin $1$. The invariant margin therefore decays like $1/\sqrt d$ in the image size
$d$ while the unconstrained margin stays constant: this is exactly the $2/\sqrt d$ collapse of
\citet{ge2021shift}, recovered as the cyclic specialization of the projection identity. The general
theorem says \emph{why} it happens: the invariance projects away every coordinate except the mean,
and a localized pattern almost entirely lives in the coordinates that get projected away.


% =====================================================================
\subsection{Q3 answered: the $\eta/L$ certificate, its missing converse, and the two-sided bracket}
\label{sub:sandwich}

\paragraph{The question this answers.}
Part~I gave the Lipschitz--margin certificate: if an invariant score $g$ has feature margin $\eta$
and Lipschitz constant $L$, then the $\ell_2$ robust radius satisfies $r_2(g,x) \ge \eta/L$. This is
the link from an invariance to an $\ell_2$ radius that no prior paper supplied. But a one-sided
certificate is a weak predictor unless we also know how much it can \emph{under}estimate the true
radius. So Q3 splits into three sub-questions. First, is the certificate ever loose by an unbounded
amount (no universal converse)? Second, under what condition can we recover an upper bound, turning
the certificate into a two-sided bracket that pins $r_2$ down? Third, when is the bracket exact?

\paragraph{Status.}
The lower bound is the textbook Lipschitz--margin certificate \citep{hein2017formal,tsuzuku2018lipschitz};
we do not claim it. The genuine contribution here is the \emph{upper arm}. \citet{tsuzuku2018lipschitz}
state a certificate chain whose top inequality (certificate $\le$ true radius) they say cannot be
guaranteed, and measure the resulting gap empirically as $1.8$--$2.4\times$; prior upper bounds on the
true radius were only empirical, read off by running a minimal-distortion attack
\citep{croce2020reliable}. Theorem~\ref{thm:sandwich} supplies that missing arm under an explicit,
checkable reachability condition (a co-Lipschitz bound along one path to the boundary). The linear
case (Corollary~\ref{cor:linear}) is the classical point-to-hyperplane distance
\citep{hein2017formal,cortes1995}, included as a tightness sanity check.

We first show the certificate has \emph{no} universal converse, so we know an upper arm must require
an extra hypothesis.

\begin{proposition}[No universal converse for $\eta/L$]
\label{prop:noconverse}
There is no universal constant $C$ such that every invariant classifier obeying the
Lipschitz--margin certificate also obeys $r_2(g,x) \le C\,\eta/L$. The failure occurs even when the
invariant feature depends only on the DC coordinate.
\end{proposition}

\begin{proof}
We exhibit a one-parameter family whose ratio $r_2/(\eta/L)$ grows without bound.
\emph{Step 1 (a DC-only invariant feature).}
Write $t = e^\top x = \fdc(x)$, the DC coordinate, which is shift-invariant. For an integer
$n \ge 1$ set the invariant scalar feature and classifier
\[
\Psi_n(x) = t^{2n+1}, \qquad g_n(x) = \Psi_n(x).
\]
Take the two data points $x_+ = e$ (label $+1$) and $x_- = -e$ (label $-1$), where $e = d^{-1/2}\one$
so that $t(x_\pm) = \pm 1$.

\emph{Step 2 (the feature margin).}
Then $g_n(e) = 1^{2n+1} = 1$ and $g_n(-e) = (-1)^{2n+1} = -1$. In the one-dimensional feature space
the two classes sit at $+1$ and $-1$, so the signed feature margin is $\eta = 1$.

\emph{Step 3 (the Lipschitz constant on the relevant segment).}
On the DC segment $\{t e : t \in [-1,1]\}$, the derivative of $\Psi_n$ in $t$ is $(2n+1)t^{2n}$, of
magnitude at most $2n+1$, attained at $t = \pm 1$. \emph{(maximum over $t$ of $|(2n+1)t^{2n}|$ on
$[-1,1]$)} Hence the Lipschitz constant is $L_n = 2n+1$, and the certificate gives
$\eta/L_n = 1/(2n+1)$.

\emph{Step 4 (the true radius).}
The decision boundary is $g_n = 0$, i.e. $t = 0$. The closest point to $x_+ = e$ with $t = 0$ is the
origin, at Euclidean distance $\|e - 0\|_2 = 1$; perturbing inside the anti-invariant subspace $\Vac$
(where $\one^\top\delta = 0$) does not change $t$ at all, so it cannot move the score. Therefore the
true robust radius is $r_2(g_n, e) = 1$.

\emph{Step 5 (the ratio is unbounded).}
Combining Steps 3 and 4,
\[
\frac{r_2(g_n, e)}{\eta/L_n} = \frac{1}{1/(2n+1)} = 2n+1 \xrightarrow[n\to\infty]{} \infty.
\]
No single constant $C$ can dominate $2n+1$ for all $n$, so no universal converse exists. Every
feature here depends only on $t$, proving the last sentence.
\end{proof}

The reason the certificate is loose in Proposition~\ref{prop:noconverse} is now visible: the feature
$t^{2n+1}$ is steep near the data ($L_n$ large) but \emph{flat} near the boundary $t=0$, so the
upper Lipschitz constant overstates how fast the score actually falls toward the boundary. Putting a
\emph{floor} under that fall, a co-Lipschitz bound (Remark~\ref{rem:terms}), is exactly what closes
the gap.

\begin{theorem}[Two-sided robust-radius bracket: a conditional converse]
\label{thm:sandwich}
Let $g$ be continuous and $L$-Lipschitz in $\ell_2$ on a convex set $S$ containing $x$, its nearest
boundary point, and the path $\gamma$ below; let $g$ be correct at $x$ with margin
$\mu = y\,g(x) > 0$, and write $\mathcal B = \{z : g(z) = 0\}$ for the decision boundary.
\begin{enumerate}[leftmargin=1.7em,label=\textup{(\roman*)}]
\item \emph{Lower arm:} $r_2(g,x) \ge \mu/L$.
\item \emph{Upper arm:} if there is a continuous path $\gamma:[0,T]\to S$ with $\gamma(0)=x$,
$g(\gamma(T))=0$, and a co-Lipschitz bound $|g(\gamma(t)) - g(x)| \ge \alpha\,\|\gamma(t)-x\|_2$ up to
the first boundary hit, then $r_2(g,x) \le \mu/\alpha$, and necessarily $\alpha \le L$.
\end{enumerate}
Hence, with the local bi-Lipschitz condition number $\kappa := L/\alpha \ge 1$,
\[
\frac{\mu}{L} \ \le\ r_2(g,x) \ \le\ \frac{\mu}{\alpha} = \kappa\cdot\frac{\mu}{L},
\]
so the certificate determines the robust radius up to the factor $\kappa$, and is exact
($r_2 = \mu/L$) if and only if $\kappa = 1$.
\end{theorem}

\begin{proof}
\emph{Lower arm.}
For any admissible $\delta$, the Lipschitz bound gives
$|g(x+\delta) - g(x)| \le L\,\|\delta\|_2$. If $\|\delta\|_2 < \mu/L$ then
$|g(x+\delta) - g(x)| < \mu = y\,g(x)$, so $y\,g(x+\delta) > 0$ and the sign cannot flip. Hence no
adversarial example exists inside radius $\mu/L$, i.e. $r_2(g,x) \ge \mu/L$. \emph{(this is the
Part~I certificate; restated for completeness)}

\emph{Upper arm.}
Evaluate the co-Lipschitz hypothesis at the first boundary hit $t = T$, where $g(\gamma(T)) = 0$:
\[
\mu = |g(x)| = |g(x) - g(\gamma(T))| \ \ge\ \alpha\,\|\gamma(T) - x\|_2,
\]
using $g(\gamma(T))=0$ in the middle and the co-Lipschitz floor on the right. Rearranging \emph{(divide
by $\alpha>0$)},
\[
\|\gamma(T) - x\|_2 \ \le\ \frac{\mu}{\alpha}.
\]
Now $\gamma(T) \in \mathcal B$ is a point where the score is zero, so an arbitrarily small further push
across $\mathcal B$ flips the prediction; the robust radius, being the \emph{distance to the boundary
set}, is at most the distance to this one boundary point: \emph{(infimum over boundary points is at
most one particular boundary point)}
\[
r_2(g,x) = \dist(x,\mathcal B) \ \le\ \|\gamma(T) - x\|_2 \ \le\ \frac{\mu}{\alpha}.
\]

\emph{The constants compare correctly.}
Along $\gamma$ the same secant $|g(\gamma(t)) - g(x)|$ is floored by $\alpha\|\gamma(t)-x\|_2$ and
capped by $L\|\gamma(t)-x\|_2$, so $\alpha\,\|\gamma(t)-x\|_2 \le L\,\|\gamma(t)-x\|_2$, hence
$\alpha \le L$ and $\kappa = L/\alpha \ge 1$. Chaining the two arms gives the bracket, and equality
$r_2 = \mu/L$ forces $\mu/\alpha = \mu/L$, i.e. $\kappa = 1$.
\end{proof}

\begin{remark}[When is the upper arm usable?]
The lower arm $r_2\ge\mu/L$ is free (it needs only the Lipschitz constant). The upper arm needs an
\emph{existence} witness: one path to the boundary with a positive co-Lipschitz slope $\alpha$. This is
not automatic on an arbitrary trained network, so the bracket is a \emph{conditional} converse, and
making it usable means producing such a path. Two cases where it is concrete: (i) for an explicit
invariant feature, such as the power-spectrum quadratic surrogate below, $\alpha$ is computable in
closed form on a ball, giving a numerical bracket; (ii) for a general differentiable score one can
descend along $-\nabla g$ from $x$ to the boundary and take $\alpha$ to be the smallest gradient norm
met along that path, which is a valid (if loose) co-Lipschitz floor. The honest gap between ``a
conditional converse exists'' and ``a tight bracket on a trained CNN'' is exactly the difficulty of
certifying a good $\alpha$, just as the lower arm's looseness is the difficulty of certifying a tight
$L$.
\end{remark}

\begin{corollary}[Linear invariant features: the bracket is exact]
\label{cor:linear}
If $g(x) = w^\top x$ with $w \in V^G$, then along the straight path $\gamma(t) = x - t\,w/\|w\|_2$ one
has $L = \alpha = \|w\|_2$, so $\kappa = 1$ and
\[
r_2(g,x) = \frac{\mu}{\|w\|_2} = \frac{\eta}{\|w\|_2},
\]
the exact point-to-hyperplane distance \citep{hein2017formal,cortes1995}.
\end{corollary}

\begin{proof}
For linear $g$, $|g(x) - g(z)| = |w^\top(x-z)| \le \|w\|_2\,\|x-z\|_2$ by Cauchy--Schwarz, with
equality exactly when $x-z$ is parallel to $w$. So the upper Lipschitz constant is $L = \|w\|_2$, and
along the straight path in direction $-w$ the inequality is an equality, giving co-Lipschitz constant
$\alpha = \|w\|_2$ as well. Thus $\kappa = 1$, and the path reaches the boundary at distance
$\mu/\|w\|_2$; Theorem~\ref{thm:sandwich} pins this as the exact radius. The margin $\mu$ equals the
feature margin $\eta$ for the max-margin separator, giving the stated form.
\end{proof}

\begin{remark}[Reconciling the no-converse counterexample with the bracket]
\label{rem:reconcile}
The bracket is not contradicted by Proposition~\ref{prop:noconverse}; it explains it. For
$\Psi_n(x) = t^{2n+1}$ at $x = e$, the straight DC path to the origin has the feature drop from $1$ to
$0$ over Euclidean distance $1$, so its secant co-Lipschitz constant is $\alpha = 1$, while the upper
Lipschitz on the segment is $L = 2n+1$. The bracket reads $[\mu/L,\,\mu/\alpha] = [1/(2n+1),\,1]$ and
correctly contains $r_2 = 1$, sitting at the $\mu/\alpha$ end. The failure of a \emph{universal}
converse is precisely that $\kappa = 2n+1$ is unbounded \emph{across the model family}, not that any
single model lacks a bracket. Note also the difference between the \emph{pointwise} lower-Lipschitz
$\inf_t|\Psi_n'|$, which vanishes at the boundary (a metric-subregularity failure
\citep{dontchev2009}) and so kills any derivative-based converse, and the \emph{secant}
co-Lipschitz above, which stays positive and yields the finite bracket.
\end{remark}

\begin{remark}[A computable instance]
For the power-spectrum surrogate of Section~\ref{sub:ps}, both arms are explicit: the upper Lipschitz
satisfies $L \le 2B$ on the ball $\|x\|_2 \le B$, and along the gradient-descent path the
co-Lipschitz constant is the minimum gradient norm $\alpha = \inf_{z\in\gamma}\|\nabla g(z)\|_2$,
strictly positive as long as the path avoids the single critical point of the quadratic. This gives
the first two-sided radius statement for an invariant-feature model, with measured condition number
$\kappa \approx 1.6$--$2.4$ and the attack-found radius lying strictly inside the bracket for every
seed.
\end{remark}


% =====================================================================
\subsection{Q2, deeper: a nonlinear invariant the linear one cannot match}
\label{sub:ps}

\paragraph{The question this answers.}
Theorem~\ref{thm:fg-margin} shows linear invariance is weak: for cyclic shifts it keeps only the
mean. The natural next step is a \emph{nonlinear} invariant feature that keeps more, while still
being computable by a standard layer. The classical choice is the power spectrum $\{|\hat x_k|^2\}_k$,
which is shift-invariant because a shift only rotates Fourier phases. The point of this subsection is
twofold: first, a single convolution--square--global-pool layer literally \emph{realizes} power-spectrum
coordinates, so this is not an exotic feature but what such a layer already computes; second, the power
spectrum is still blind to \emph{phase}, and a degree-three invariant, the bispectrum, separates
phase-coded classes that the power spectrum cannot, including DC-matched pairs that no shift-invariant linear feature separates either.

\paragraph{Status.}
Both facts are known-and-reproved. Power spectra, autocorrelation, and scattering invariants are
classical \citep{brunamallat2012}; the realization lemma is a low-level architectural identity (with
a corrected DFT normalization and real-filter handling). The bispectrum and its completeness-up-to-shift
are classical \citep{invariants_kakarala}. The only fresh angle is the \emph{application}: using a
higher-order invariant to defeat exactly the localized, phase-coded corner where Ge's forced trade-off
bites, and quantifying its cost as a larger Lipschitz constant, hence a smaller $\eta/L$ at equal margin.

\begin{lemma}[A conv--square--pool layer realizes the power spectrum]
\label{lem:ps}
Use the unitary DFT $\hat x_k = d^{-1/2}\sum_j x_j e^{-2\pi i jk/d}$. Define the unnormalised circular
cross-correlation $(w \star x)_s = \sum_j w_j x_{j+s \bmod d}$ and the feature
\[
\Phi_w(x) := \frac1d \sum_{s=0}^{d-1} (w\star x)_s^2 ,
\]
i.e. correlate $x$ with a filter $w$, square the response, and global-average-pool. Then for real
$w,x$,
\[
\Phi_w(x) = \sum_{k=0}^{d-1} |\hat w_k|^2\,|\hat x_k|^2 ,
\]
a nonnegative-weighted combination of the power-spectrum coordinates $|\hat x_k|^2$. Real filters can
isolate every nonredundant real power coordinate ($|\hat x_0|^2$, $|\hat x_{d/2}|^2$ when $d$ is even,
and each conjugate-pair sum $|\hat x_k|^2 + |\hat x_{d-k}|^2$).
\end{lemma}

\begin{proof}
\emph{Step 1 (DFT of a correlation).}
By the convolution/correlation theorem for the unitary DFT, the transform of $w\star x$ is
$\widehat{w\star x}_k = \sqrt{d}\,\overline{\hat w_k}\,\hat x_k$ \emph{(the $\sqrt d$ is the unitary
normalization; the conjugate is because correlation flips one argument)}.

\emph{Step 2 (Parseval turns the pooled square into a frequency sum).}
The pooling sum is a squared $\ell_2$ norm in space, and the unitary DFT preserves $\ell_2$ norms
(Parseval):
\[
\Phi_w(x) = \frac1d\sum_s |(w\star x)_s|^2 = \frac1d\sum_k |\widehat{w\star x}_k|^2 .
\]

\emph{Step 3 (substitute Step 1).}
Insert $\widehat{w\star x}_k = \sqrt d\,\overline{\hat w_k}\hat x_k$ and take squared modulus:
\[
\Phi_w(x) = \frac1d\sum_k \big|\sqrt d\,\overline{\hat w_k}\hat x_k\big|^2
          = \frac1d\sum_k d\,|\hat w_k|^2|\hat x_k|^2
          = \sum_k |\hat w_k|^2\,|\hat x_k|^2 ,
\]
the $\sqrt d$ squaring to $d$ and cancelling the $1/d$. This is the claimed identity.

\emph{Step 4 (which coordinates real filters reach).}
For real $w$, conjugate symmetry $\hat w_{d-k} = \overline{\hat w_k}$ forces $|\hat w_{d-k}|^2 =
|\hat w_k|^2$, so a real filter weights $|\hat x_k|^2$ and $|\hat x_{d-k}|^2$ \emph{equally}; it cannot
separate a conjugate pair. But it \emph{can} isolate a pair: choose $\hat w$ supported only on
$\{k, d-k\}$ with conjugate values (which makes $w$ real), giving exactly the pair sum
$|\hat x_k|^2 + |\hat x_{d-k}|^2$; the self-conjugate coordinates $k=0$ and (for even $d$) $k=d/2$ are
isolated singly. Linear combinations of these give the stated real span.
\end{proof}

So one standard layer escapes the one-dimensional invariant linear family of
Theorem~\ref{thm:fg-margin} and sees the whole power spectrum. But the power spectrum still discards
phase, and that is exactly what Ge's localized example exploits.

\begin{proposition}[The bispectrum separates phase-coded classes the power spectrum cannot]
\label{prop:bispectrum}
The bispectrum $B_{k_1,k_2}(x) = \hat x_{k_1}\hat x_{k_2}\overline{\hat x_{k_1+k_2}}$ is
shift-invariant and is realizable as a degree-three convolution--product--pool feature. There exist
classes with \emph{identical} power spectra that a single bispectrum coordinate separates: for Ge's
dot, $x_+ = e_j$ and $x_- = -e_j$ have $|\hat x_k|^2 = 1/d$ for all $k$, so every power-spectrum
feature is blind to the difference, yet
\[
B_{0,0}(e_j) = +d^{-3/2}, \qquad B_{0,0}(-e_j) = -d^{-3/2},
\]
separating them with margin $d^{-3/2}$. (The linear DC invariant also separates this particular pair,
through $\fdc(\pm e_j) = \pm d^{-1/2}$; the sharper claim, proved in Step~3, is that the bispectrum
separates \emph{any} two non-shift-equivalent orbits, including pairs with identical DC and identical
power spectrum, which no shift-invariant linear feature can do.) This separability comes at the cost
of a larger Lipschitz constant, $O(B^2)$ versus $O(B)$ for the power spectrum on the ball
$\|x\|_2 \le B$, i.e. a smaller certified $\eta/L$ at equal margin.
\end{proposition}

\begin{proof}
\emph{Step 1 (shift-invariance of the bispectrum).}
Under $x \mapsto S_s x$ each Fourier coordinate picks up a unit phase, $\hat x_k \mapsto \omega^{sk}\hat
x_k$ with $\omega = e^{-2\pi i/d}$ (the forward DFT convention of Part~I; the shift theorem
$\mathrm{DFT}(R_\tau x)_k = \omega^{k\tau}\hat x_k$). The three factors of $B_{k_1,k_2}$ acquire phases $\omega^{s k_1}$,
$\omega^{s k_2}$, and (from the conjugate) $\omega^{-s(k_1+k_2)}$, whose product is
$\omega^{s(k_1 + k_2 - k_1 - k_2)} = \omega^0 = 1$. The phases cancel, so $B_{k_1,k_2}$ is invariant.
It is a cubic form in $x$, hence realizable by circular correlations, a pointwise product, and global
average pooling (a degree-three analogue of Lemma~\ref{lem:ps}).

\emph{Step 2 (the dot has a flat power spectrum).}
For $x = e_j$ with the unitary DFT, $\hat x_k = d^{-1/2}\omega^{jk}$ where $\omega = e^{-2\pi i/d}$, so
$|\hat x_k|^2 = 1/d$ for every $k$ (the phase has modulus one). The same holds for $-e_j$. The two
power spectra are identical, so by Lemma~\ref{lem:ps} no power-spectrum feature distinguishes them.
A shift-invariant \emph{linear} feature sees only the DC value (Theorem~\ref{thm:fg-margin}); for this
particular pair the DC value $\fdc(\pm e_j) = \pm d^{-1/2}$ does differ, so the linear invariant
separates $\pm e_j$, but only with the weak $1/\sqrt d$ margin of the Ge collapse, and it would be
\emph{completely} blind to any DC-matched pair such as two shifts of a phase-coded signal.

\emph{Step 3 (a bispectrum coordinate separates them).}
Take $k_1 = k_2 = 0$. Then $B_{0,0}(x) = \hat x_0 \hat x_0 \overline{\hat x_0} = \hat x_0\,|\hat x_0|^2$.
For $e_j$, $\hat x_0 = d^{-1/2}$ and $|\hat x_0|^2 = 1/d$, so $B_{0,0}(e_j) = d^{-1/2}\cdot d^{-1} =
d^{-3/2}$. For $-e_j$ the sign of $\hat x_0$ flips, giving $-d^{-3/2}$. The single coordinate
$B_{0,0}$ takes opposite signs on the two classes, separating them with margin $d^{-3/2}$. More
generally, when all Fourier coefficients are nonzero the bispectrum determines the signal up to a
global shift \citep{invariants_kakarala}, so it separates any two shift-orbits that are not
shift-equivalent.

\emph{Step 4 (the Lipschitz cost).}
A cubic form has gradient of order $\|x\|_2^2$, so on the ball $\|x\|_2 \le B$ the bispectrum is
locally Lipschitz with constant $O(B^2)$, versus $O(B)$ for the quadratic power spectrum. At equal
margin a larger $L$ means a smaller certified $\eta/L$, which is the stated trade: the higher-order
invariant buys separability of the phase-coded corner at the price of a smaller certificate.
\end{proof}


% =====================================================================
\subsection{Q5 answered, part one: on orbit-closed data invariance is margin-free}
\label{sub:selfsym}

\paragraph{The question this answers.}
The reproduction's quadrant diagnostic (Part~II, Act~II) plotted models on a
consistency-versus-robustness plane and could not say \emph{why} a model lands where it lands, nor
whether building in invariance helps. Q5 asks for
the optimization story. The first half of the answer is a surprising \emph{negative} result: when the
data is orbit-closed (every shift of a training point is also a training point with the same label),
imposing exact shift-invariance does not change the max-margin value at all. Gradient descent on the
unconstrained network already self-symmetrizes for free, so there is no margin to be gained by tying
weights. This explains the experiment in which unconstrained and shift-tied networks reach essentially
the same robustness on generic orbit data.

\paragraph{Status.}
A medium-strength novel candidate. It is the nonlinear-ReLU, orbit-closed analogue of the linear
steerable/augmentation equivalence of \citet{chenzhu2023}; the linear case is theirs, the ReLU
min-norm accounting below is the new content. We use one architectural identity from Part~I: a
convolution-plus-global-pool unit is the same function as a two-layer net whose hidden weights are
tied across the shift orbit of one filter (the ``orbit bundle''). We also use the standard fact that a
positively homogeneous ReLU unit $a\,\sigma(w^\top x)$ can be rescaled $(a,w)\mapsto(a/c,\,cw)$ without
changing its function, so its minimum-norm representation costs $|a|\,\|w\|$ (taking $c^2 = |a|/\|w\|$
balances $a^2$ and $\|w\|^2$).

\begin{theorem}[Invariance is margin-free on orbit-closed data]
\label{thm:selfsym}
Let the data be orbit-closed with labels constant on orbits. For two-layer ReLU networks
$f(x) = \sum_r a_r \sigma(w_r^\top x)$ (positively homogeneous, no bias), the minimum parameter norm
$\tfrac12\sum_r(a_r^2 + \|w_r\|^2)$ achieving unit margin is the \emph{same} for the unconstrained
class and for the shift-invariant (tied orbit-bundle) class. Equivalently, imposing exact
shift-invariance changes neither the max-margin value nor the achievable margin.
\end{theorem}

\begin{proof}
Write $C^\star_{\mathrm{unc}}$ and $C^\star_{\mathrm{inv}}$ for the two minimum costs. We show
$C^\star_{\mathrm{inv}} \ge C^\star_{\mathrm{unc}}$ and $C^\star_{\mathrm{inv}} \le
C^\star_{\mathrm{unc}}$.

\emph{Step 1 ($\ge$, the easy direction).}
The tied (invariant) class is a \emph{subset} of the unconstrained class: every invariant network is
in particular some network. Minimizing the same cost over a smaller feasible set cannot give a smaller
value, so $C^\star_{\mathrm{inv}} \ge C^\star_{\mathrm{unc}}$.

\emph{Step 2 (the symmetrization map).}
Let $f^\star = \sum_r a_r\sigma(w_r^\top x)$ be any unconstrained network with margin $\ge 1$. Define
its orbit average
\[
f_{\mathrm{sym}}(x) = \frac1d\sum_{s=0}^{d-1} f^\star(S_s x).
\]
Reindexing the pooling sum shows $f_{\mathrm{sym}}$ is shift-invariant, and by the orbit-bundle
identity it lies in the tied class.

\emph{Step 3 (symmetrization preserves the margin).}
Fix a training point $(x_i, y_i)$. Because the data is orbit-closed with orbit-constant labels, each
shifted point $S_s x_i$ is itself a training point with the same label $y_i$, so $y_i f^\star(S_s x_i)
\ge 1$ for every $s$. Averaging over $s$ \emph{(average of numbers each $\ge 1$ is $\ge 1$)},
\[
y_i f_{\mathrm{sym}}(x_i) = \frac1d\sum_s y_i f^\star(S_s x_i) \ge \frac1d\sum_s 1 = 1 .
\]
So $f_{\mathrm{sym}}$ still has unit margin.

\emph{Step 4 (the min-norm cost of one unit).}
For a single ReLU unit $a\,\sigma(w^\top x)$, positive homogeneity of $\sigma$ means
$a\,\sigma(w^\top x) = (a/c)\,\sigma((cw)^\top x)$ for any $c>0$, so we may rescale freely. The cost
$\tfrac12((a/c)^2 + \|cw\|^2)$ is minimized over $c$ by the arithmetic--geometric-mean balance at
$c^2 = |a|/\|w\|$, giving minimal cost $|a|\,\|w\|$. \emph{(this is the standard min-norm value of one
homogeneous unit)}

\emph{Step 5 (the bundle realizing $f_{\mathrm{sym}}$ costs no more than $f^\star$).}
Expanding the average, $f_{\mathrm{sym}}(x) = \frac1d\sum_r\sum_s a_r\,\sigma(w_r^\top S_s x)
= \frac1d\sum_r\sum_s a_r\,\sigma((S_{-s}w_r)^\top x)$, using that $S_s$ is orthogonal so
$w_r^\top S_s x = (S_s^\top w_r)^\top x = (S_{-s}w_r)^\top x$. This is a network whose units are
$\{(a_r/d,\ S_{-s}w_r)\}_{r,s}$. By Step~4 its min-norm cost is
\[
\sum_r\sum_s \frac{|a_r|}{d}\,\|S_{-s}w_r\|
= \sum_r\sum_s \frac{|a_r|}{d}\,\|w_r\|
= \sum_r |a_r|\,\|w_r\| ,
\]
where the shift $S_{-s}$ preserves the norm (orthogonality) and the inner sum over the $d$ shifts of
the factor $1/d$ collapses. But $\sum_r |a_r|\,\|w_r\|$ is exactly the min-norm cost of the original
$f^\star$ by Step~4. Hence $f_{\mathrm{sym}}$ is realized in the tied class at cost no larger than the
cost of $f^\star$.

\emph{Step 6 (conclude).}
Taking $f^\star$ to be a unconstrained \emph{minimizer}, Step~5 produces a feasible invariant network
of the same cost, so $C^\star_{\mathrm{inv}} \le C^\star_{\mathrm{unc}}$. With Step~1, the two are
equal. There is therefore an invariant max-margin solution of the same norm: invariance is
margin-free here.
\end{proof}

\begin{remark}[What it does and does not give]
\label{rem:selfsym}
Theorem~\ref{thm:selfsym} equalizes the \emph{margin}; it does not say which max-margin solution
gradient descent selects, and the certified radius $\eta/L$ depends on the realized Lipschitz term
$L$, not the margin alone. So invariance can still change $\eta/L$ by changing $L$ even when it cannot
change $\eta$. That residual selection question, on data where the orbit-closed premise fails, is the
content of the next subsection and its open conjecture.
\end{remark}


% =====================================================================
\subsection{Q5 answered, part two: in the lazy regime invariance lowers the Lipschitz constant}
\label{sub:lazy}

\paragraph{The question this answers.}
Theorem~\ref{thm:selfsym} says invariance cannot help \emph{the margin} on orbit-closed data. But on
data with near-orthogonal cluster geometry, experiment shows it \emph{does} help $\eta/L$, and the
gap is largest in the lazy (NTK) training regime. The remaining question is whether that advantage is
provable. The answer is yes in the lazy, smooth-activation regime: there the trained network is the
minimum-RKHS-norm interpolant of a fixed kernel, orbit-averaging is an orthogonal projection in that
RKHS, and a projection cannot increase the norm, hence cannot increase the Lipschitz constant, at
equal margin. So invariance buys robustness through a smaller $L$, exactly the lever
Remark~\ref{rem:selfsym} left open.

\paragraph{Status.}
A medium-strength novel candidate, with an honestly limited scope. It proves the cluster conjecture
\emph{only} in the lazy, smooth-activation regime, by combining the lazy/NTK fixed-kernel picture
\citep{jacot2018,chizat2019} with the RKHS orbit-averaging projection of \citet{elesedy2021} and
the invariant-kernel symmetrization of \citet{haasdonk2007}. The smooth-activation caveat is
real and we state it: the bare two-layer ReLU NTK feature map is \emph{not} Lipschitz
\citep{biettimairal2019}, so the Lipschitz step needs a smooth activation; the rich (small-init)
and bare-ReLU trajectories stay open. From Part~I we use the lazy regime (a wide network behaves like
a fixed kernel), the RKHS norm and its Lipschitz bound $|f(x)-f(y)| \le \|f\|_K\,L_\Phi\,\|x-y\|_2$,
and orthogonal projection (a projection is non-expansive and gives a Pythagorean norm split).

\begin{theorem}[Invariance lowers the Lipschitz constant in the lazy, smooth-activation regime]
\label{thm:lazy}
Work in the lazy/NTK regime, where each trained network equals the minimum-RKHS-norm interpolant of
its fixed neural tangent kernel $K$ \citep{jacot2018,chizat2019}, with a smooth activation so that the
tangent feature map $\Phi$ is $L_\Phi$-Lipschitz on the data ball \citep{biettimairal2019}. Let the
data be orbit-closed with orbit-constant labels and let $K$ be $G$-invariant, so the Reynolds average
$\Pi_G f = \frac1{|G|}\sum_{g\in G} f(g\cdot\,)$ is defined on the RKHS $\mathcal H_K$ and the
shift-tied bundle realizes exactly the invariant subspace $\mathcal H^G = \Pi_G\mathcal H_K$ with the
restricted norm \citep{haasdonk2007}. Let $f_{\mathrm{unc}}$ be the unconstrained min-norm
interpolant and $f_{\mathrm{inv}}$ the min-norm interpolant in $\mathcal H^G$. Then
\begin{enumerate}[leftmargin=1.7em,label=\textup{(\roman*)}]
\item $\Pi_G$ is the orthogonal projection of $\mathcal H_K$ onto $\mathcal H^G$, non-expansive, with
$\|f\|_K^2 = \|\Pi_G f\|_K^2 + \|(I-\Pi_G)f\|_K^2$ \citep{elesedy2021};
\item $\Pi_G f_{\mathrm{unc}}$ still interpolates with the same unit margin, so
$\|f_{\mathrm{inv}}\|_K \le \|\Pi_G f_{\mathrm{unc}}\|_K \le \|f_{\mathrm{unc}}\|_K$;
\item hence $L_{\mathrm{inv}} \le L_\Phi\,\|f_{\mathrm{inv}}\|_K \le L_\Phi\,\|f_{\mathrm{unc}}\|_K$,
and at the common unit margin the invariant model has the weakly larger certified ratio $\eta/L$. The
gap is the discarded non-invariant energy $\|(I-\Pi_G)f_{\mathrm{unc}}\|_K^2$ (the orbit-variance),
strict unless $f_{\mathrm{unc}}$ is already invariant.
\end{enumerate}
\end{theorem}

\begin{proof}
\emph{(i) $\Pi_G$ is an orthogonal projection.}
Because $K$ is $G$-invariant, the averaging operator $\Pi_G$ maps $\mathcal H_K$ into itself, and it
is idempotent ($\Pi_G^2 = \Pi_G$, the same group-reindexing as in Theorem~\ref{thm:fg-margin}, Step~2)
and self-adjoint on $\mathcal H_K$. An idempotent self-adjoint operator has spectrum in $\{0,1\}$ and
is the orthogonal projection onto its range, which is the invariant subspace $\mathcal H^G$. For an
orthogonal projection, every $f$ splits as $f = \Pi_G f + (I-\Pi_G)f$ into orthogonal pieces, so by
the Pythagorean theorem $\|f\|_K^2 = \|\Pi_G f\|_K^2 + \|(I-\Pi_G)f\|_K^2$
\citep{elesedy2021}; dropping the second nonnegative term gives non-expansiveness
$\|\Pi_G f\|_K \le \|f\|_K$.

\emph{(ii) the projected interpolant is still feasible, with no larger norm.}
Orbit-constant labels mean $y_i f_{\mathrm{unc}}(g\cdot x_i) \ge 1$ for every $g$ (each $g\cdot x_i$
is a training point of the same label, interpolated by $f_{\mathrm{unc}}$). Averaging over $g$,
\[
y_i\,(\Pi_G f_{\mathrm{unc}})(x_i) = \frac1{|G|}\sum_g y_i f_{\mathrm{unc}}(g\cdot x_i) \ge 1,
\]
so $\Pi_G f_{\mathrm{unc}} \in \mathcal H^G$ is feasible for the invariant min-norm problem. Since
$f_{\mathrm{inv}}$ is the \emph{minimizer} of that problem, $\|f_{\mathrm{inv}}\|_K \le
\|\Pi_G f_{\mathrm{unc}}\|_K$, and by (i) this is $\le \|f_{\mathrm{unc}}\|_K$.

\emph{(iii) smaller norm gives smaller Lipschitz, hence larger $\eta/L$.}
For any RKHS function, the reproducing property and Cauchy--Schwarz give
\[
|f(x) - f(y)| = |\langle f,\ \Phi(x) - \Phi(y)\rangle_K|
\le \|f\|_K\,\|\Phi(x) - \Phi(y)\|_K
\le \|f\|_K\,L_\Phi\,\|x - y\|_2,
\]
the last step using that $\Phi$ is $L_\Phi$-Lipschitz on the data ball (smooth activation)
\citep{biettimairal2019}. So a valid Lipschitz constant is $L \le L_\Phi\,\|f\|_K$. Applying this to
each interpolant and inserting (ii), $L_{\mathrm{inv}} \le L_\Phi\,\|f_{\mathrm{inv}}\|_K \le
L_\Phi\,\|f_{\mathrm{unc}}\|_K$. The two models share the same unit margin $\eta$, so the invariant
model's certified $\eta/L_{\mathrm{inv}}$ is at least the unconstrained model's. The size of the gap is
controlled by the discarded energy $\|(I-\Pi_G)f_{\mathrm{unc}}\|_K^2$ from the split in (i): it is
zero exactly when $f_{\mathrm{unc}}$ is already invariant, and strictly positive otherwise.
\end{proof}

\begin{remark}[The ReLU caveat, stated honestly]
\label{rem:relu}
Steps (i)--(ii) are activation-independent: they are pure RKHS geometry. Only step (iii) uses
smoothness. The bare two-layer ReLU NTK feature map is \emph{not} Lipschitz; its RKHS contains
unit-norm functions of unbounded Lipschitz constant \citep{biettimairal2019}, so for ReLU one must
either take a smooth activation or carry $\sup\|\nabla\Phi\|$ on the data ball explicitly. Within that
scope the theorem proves the cluster conjecture, exactly in the regime where the project's
experiments report the largest gap (orbit-variance $\approx 0.3$, tied-over-unconstrained $\eta/L$ up
to $4.6\times$). The rich (small-init) regime and the bare-ReLU trajectory remain open.
\end{remark}


% =====================================================================
\subsection{The open frontier: the rich-regime cluster conjecture and how to attack it}
\label{sub:conjecture}

Theorem~\ref{thm:lazy} closes the lazy half. The rich (feature-learning, small-init) ReLU half is
genuinely open. We state it as a conjecture, then, in the spirit of a primer, list the candidate
proof paths and the Part~I tools each one needs, so a collaborator can pick one up.

\begin{conjecture}[Invariance helps on cluster geometry]
\label{conj:cluster}
On near-orthogonal cluster-geometry orbit data, the tied-invariant two-layer ReLU network's gradient
flow reaches a strictly larger $\eta/L$ than the unconstrained network's, with the gap growing as
initialization moves toward the lazy regime. Equivalently, weight sharing plus pooling re-points the
implicit bias toward the robust invariant solution.
\end{conjecture}

Why it is hard, in one line each: the orbit data is rank-structured, \emph{not} the near-orthogonal
$\Omega(d)$-cluster model whose non-robust selection is known, so that conclusion does not transfer;
the tied program has its own constrained KKT conditions that must be written, not inherited;
homogeneity gives convergence to \emph{some} KKT point but not \emph{which} one; and $\eta/L$ depends on
the realized Lipschitz term, so one must track the input-gradient norm along the trajectory, not just
the functional margin.

\paragraph{Candidate proof paths and the tools each needs.}
\begin{enumerate}[leftmargin=1.6em]
\item \textbf{Frequency-domain reduction (one orbit per class).} Replace the ReLU head by the
conv--square--pool surrogate so the features are power-spectrum coordinates (Lemma~\ref{lem:ps}); the
invariant margin program then becomes a linear/quadratic max-margin in Fourier-magnitude space,
possibly solvable in closed form, and one compares its $\eta/L$ to the unconstrained KKT point on the
same data, then swaps the surrogate for ReLU.
\emph{Tools:} the \emph{Fourier/power-spectrum invariants} and the \emph{convex-hull max-margin}
geometry, both introduced in the invariance and learning clusters of Part~I.
\item \textbf{Explicit KKT in Fourier coordinates.} Use the circulant/DFT diagonalization of the shift
action so that frequencies decouple, write the unconstrained and constrained KKT systems per
frequency, and solve frequency by frequency.
\emph{Tools:} the \emph{implicit-bias/KKT machinery} for homogeneous networks
\citep{lyuli2020,jitelgarsky2019,soudry2018} from the learning cluster of Part~I, plus the
\emph{per-irreducible (DFT) decomposition} of the shift group from the invariance cluster.
\item \textbf{Rich versus lazy contrast.} Run the same construction at small init (rich,
feature-learning) and at NTK scale (lazy), and show the selected $\eta/L$ differs, with the invariant
constraint mattering most in the rich regime.
\emph{Tools:} the \emph{NTK/lazy-versus-rich} distinction from the learning cluster of Part~I; the
lazy end is already pinned by Theorem~\ref{thm:lazy}.
\item \textbf{Bound the two ends separately.} Lower-bound the invariant network's achieved margin by
the projected convex-hull margin of Theorem~\ref{thm:fg-margin} in feature space, upper-bound the
unconstrained network's $\eta/L$ via a feature-averaging direction, then compare.
\emph{Tools:} the \emph{projection-margin accounting} of Section~\ref{sub:fg-margin} and the
\emph{feature-averaging implicit-bias} results \citep{li2025feature} from the learning cluster.
\end{enumerate}
Each route also needs the supporting setup lemmas, the bare $\eta/L$ lower bound, the orbit-rank and
orbit-bundle identities, and the KKT scaffolding; we have stated here only the load-bearing results
and leave the routine scaffolding to the reader who picks up a route.


% =====================================================================
\subsection{Tying theory to the benchmark: the capacity-matched $\eta/L$ decomposition}
\label{sub:empirical}

\paragraph{The question this answers.}
Q1 and Q4: if shift-consistency does not order robustness, what single measurable quantity does, and
does the answer survive a strong attack at matched capacity? The empirical headline is that $\eta/L$,
computed in the invariant feature space, is that quantity, and the theory above predicts \emph{which}
of its two parts a given architecture moves.

\paragraph{Status.}
This is the best empirical candidate, and its framing is the honest one of
Remark~\ref{rem:calibration}: not ``does invariance help?'' but ``when capacity is controlled, which
term, margin $\eta$ or Lipschitz $L$, does a given invariance move?'' Robustness is measured by a
standardized strong attack (AutoAttack) at matched parameter counts \citep{croce2020reliable,robustbench2020},
with capacity controlled as in architecture-robustness studies \citep{robarch2023}.

\paragraph{The findings, briefly.}
Across capacity-matched arms (standard, anti-aliased, exact-cyclic, shift-augmented, identical
parameter counts), $\eta/L$ versus the $\ell_2$ robust radius tracks at Pearson $0.998$, while
shift-consistency \emph{anti}-predicts robustness. The decomposition attributes mechanism exactly as
the theory says it should: anti-aliasing raises $\eta/L$ by lowering the Lipschitz term $L$
(Theorem~\ref{thm:lazy}'s lever), and exact cyclic invariance lowers $\eta/L$ by shrinking the margin
$\eta$ (the real-data form of the Ge collapse of Section~\ref{sub:fg-margin}). Under adversarial
training, where AutoAttack becomes informative with no gradient masking, shift-consistency still
anti-predicts robustness (Pearson $-0.88$), and the \emph{threat-matched} ratio
$\eta/\|\nabla M\|_1$, using the $\ell_\infty$ dual norm, predicts AutoAttack robust accuracy at
Pearson $0.90$, whereas the mismatched $\ell_2$ ratio does not. The two laws replicate on MNIST and
Fashion-MNIST and under an $\ell_2$ threat. This is the validated margin/Lipschitz decomposition that
turns the contradictory ``helps or hurts'' literature into a single sharp question with a measurable
answer.


% =====================================================================
% CITATION KEYS USED (all already present in primer.bib unless flagged NEW):
%   bennett2000          (existing) — convex-hull / SVM duality, Lemma hull
%   burges1998           (existing) — SVM tutorial, Lemma hull
%   cortes1995           (existing) — SVM / point-to-hyperplane, Cor linear
%   hein2017formal       (existing) — Lipschitz-margin / exact linear radius
%   tsuzuku2018lipschitz (existing) — certificate chain whose upper arm is completed
%   croce2020reliable    (existing) — empirical (attack) upper bound replaced
%   soudry2018           (existing) — implicit bias, linear separable
%   lyuli2020            (existing) — homogeneous-net KKT margin
%   jitelgarsky2019      (existing) — implicit bias / KKT
%   jacot2018            (existing) — NTK / lazy
%   chizat2019           (existing) — lazy training
%   li2025feature        (existing) — feature-averaging implicit bias
%   invariants_kakarala  (existing) — bispectrum completeness
%
%   ---- NEW KEYS (NOT yet in primer.bib — must be added) ----
%   ge2021shift               NEW — Ge et al., Shift Invariance Can Reduce Adversarial Robustness, NeurIPS 2021
%                              (the base result; cyclic 2/sqrt(d) collapse, projected-margin special case)
%   brunamallat2012      NEW — Bruna & Mallat, Invariant Scattering Convolution Networks, 2012/2013
%                              (classical translation-invariant nonlinear features)
%   chenzhu2023          NEW — Chen & Zhu, implicit bias of equivariant steerable networks, 2023
%                              (linear steerable/augmentation equivalence extended by thm:selfsym)
%   elesedy2021    NEW — Elesedy, kernel/RKHS orbit-averaging projection lemma, 2021
%                              (orthogonal-projection + Pythagorean split for thm:lazy)
%   haasdonk2007 NEW — Haasdonk & Burkhardt, invariant kernels by symmetrization, 2007
%                              (invariant tangent kernel = symmetrized kernel)
%   biettimairal2019  NEW — Bietti & Mairal, group invariance / smoothness of (convolutional) NTK, 2019
%                              (smooth-activation Lipschitz; bare-ReLU NTK non-Lipschitz caveat)
%   jacot2018, chizat2019 are the existing lazy/NTK keys reused here; the TR uses Chizat2018 but the
%   primer scheme is chizat2019 (matched to primer.bib).
%   dontchev2009 NEW — Dontchev & Rockafellar, metric (sub)regularity, 2009
%                              (cited once in rem:reconcile for the subregularity-failure aside; OPTIONAL —
%                               drop the parenthetical if the key is not added)
%   croce2020reliable       NEW — Croce & Hein, AutoAttack, 2020 (standardized strong attack)
%                              NOTE: distinct from existing croce2020reliable; if they are the same entry
%                              in primer.bib, replace croce2020reliable with croce2020reliable.
%   robustbench2020      NEW — Croce et al., RobustBench, 2020/2021 (standardized robustness benchmark)
%   robarch2023          NEW — Peng et al. (RobArch), capacity-controlled architecture robustness, 2023
% =====================================================================
